\documentclass[../main.tex]{subfiles}
\begin{document}

\section{Appendix}
\label{sect:appendix}


\subsection{Assignment 2 }
\label{apx:assignment_2_}


\subsubsection{GDT Implementation}
\label{apx:gdt_implementation}


\begin{listing}[H]
    \begin{minted}[
        linenos,
        breaklines,
        frame=single,
        framesep=3mm,
        baselinestretch=1.1,
        numbersep=6pt,
        fontsize=\footnotesize,
        tabsize=4,
        bgcolor=white,
        style=friendly,
    ]{c}
#ifndef GDT_H
#define GDT_H

#include "libc/stdint.h"

// Defines the functions and structs used for initialising a GDT in our
// operating system.

// One entry in the GDT (8 bytes).
//
// The variables defined in the GDT table might be spread, therefore you will
// see a high and low section for the variables limit and base.
typedef struct gdt_entry {
  uint16_t limit_low;  // Limit is the size of the corresponding segment
  uint16_t base_low;   // Base is the starting memory address of the segment.
  uint8_t base_middle; //
  uint8_t access;      // Access byte, hold permissions settings, such as DPL
  uint8_t limit_and_flags; // Upper 4 bits: flags , Lower 4 bits: limit
  uint8_t base_high;
} __attribute__((
    packed)) gdt_entry_t; // packed ensures compiler does not alter the
                          // memory layout of the gdt_entry struct.

// The GDT pointer, this is loaded with the lgdt assembly instruction.
typedef struct {
  uint16_t limit; // Size of GDT in bytes, subtracted by 1
  uint32_t base;  // Memory address to the start of the GDT
} __attribute__((packed)) gdt_ptr_t;

// See the gdt.c file for more information.
void gdt_init(void);

// Function defined in the multiboot assembly file which loads the
// GDT pointer
extern void gdt_flush(uint32_t gdt_ptr_address);

#endif
    \end{minted}
    \caption{Implementation of \texttt{gdt.h}}
    \label{lst:gdt_h}
\end{listing}


\begin{listing}[H]
    \begin{minted}[
        linenos,
        breaklines,
        frame=single,
        framesep=3mm,
        baselinestretch=1.1,
        numbersep=6pt,
        fontsize=\footnotesize,
        tabsize=4,
        bgcolor=white,
        style=friendly,
    ]{c}
#include "../../include/gdt.h"

// Defines GDT array and pointer to the array
static gdt_entry_t gdt[3];
static gdt_ptr_t gdt_pointer;

// Helper function to populate global descriptor table, since the memory layout
// is akward due to historical reasons.
static void gdt_set_entry(int index, uint32_t base, uint32_t limit,
                          uint8_t access, uint8_t flags) {
  gdt[index].limit_low = limit & 0xFFFF;
  gdt[index].base_low = base & 0xFFFF;
  gdt[index].base_middle = (base >> 16) & 0xFF;
  gdt[index].access = access;
  gdt[index].limit_and_flags = (flags << 4) | ((limit >> 16) & 0x0F);
  gdt[index].base_high = (base >> 24) & 0xFF;
}

// Initialises up the global descriptor table and then loads it into the CPU
void gdt_init(void) {
  // The first GDT entry should always be null
  gdt_set_entry(0, 0, 0, 0, 0);

  // Defines a code segment with access to the entire memory range.
  // access 0x9A,
  // What each bit in the access byte defines can be found in
  // https://wiki.osdev.org/Global_Descriptor_Table
  //
  // The access byte represents the following:
  // Bit 7 present bit = 1, the segment is valid.
  // Bits 6-5 DPL = 00, sets segment to highest privilege.
  // Bit 4 descriptor type = 1, defines a code or datasegment
  // Bit 3 Executable bit = 1 for a code segment.
  // Bit 2 Conforming bit = 0, Code in the segment can only be exectued from the
  // same DPL level.
  // Bit 1 Read = 1, The segment is readable.
  // Bit 0 Accessed bit = 0. The CPU will set this when it uses the segment.
  //
  // flags 0xC
  // Bit 3 = 1 Page granularity, instead of byte granularity.
  // Bit 2 = 1, defines a 32-bit protected mode segment
  // Bit 1 = 0, means the segment is NOT a 64 bit code segment
  // Bit 0 reserved bit,
  // The flags represent:
  gdt_set_entry(1, 0, 0xFFFFF, 0x9A, 0xC);

  // Defines a data segment with access to the entire memory range.
  // access 0x92,
  // The access byte represents the following:
  // Bit 7 present bit = 1, the segment is valid.
  // Bits 6-5 DPL = 00, sets segment to highest privilege.
  // Bit 4 descriptor type = 1, defines a code or datasegment
  // Bit 3 Executable bit = 0 for a data segment.
  // Bit 2 Conforming bit = 0, Code in the segment can only be exectued from the
  // same DPL level.
  // Bit 1 Write = 1, The segment is writable
  // Bit 0  Accessed bit = 0. The CPU will set this when it uses the segment.
  //
  // flags 0xC
  // Same flags as previous entry
  gdt_set_entry(2, 0, 0xFFFFF, 0x92, 0xC);

  // Defines the limit and base of our gdt pointer
  gdt_pointer.limit = sizeof(gdt) - 1;
  gdt_pointer.base = (uint32_t)&gdt;

  // Defined in the multiboot2.asm file
  gdt_flush((uint32_t)&gdt_pointer);
}
    \end{minted}
    \caption{Implementation of \texttt{gdt.c}}
    \label{lst:gdt_c}
\end{listing}


\begin{listing}[H]
    \begin{minted}[
        linenos,
        breaklines,
        frame=single,
        framesep=3mm,
        baselinestretch=1.1,
        numbersep=6pt,
        fontsize=\footnotesize,
        tabsize=4,
        bgcolor=white,
        style=friendly,
    ]{nasm}
extern main
global _start
global gdt_flush  ; makes gdt_flush visible to C code
global idt_flush  ; makes idt_flush visible to C code

section .multiboot_header
header_start:
    dd 0xe85250d6 	                                                ; Magic number (multiboot 2)
    dd 0				                                            ; Architecture 0 (protected mode i386)
    dd header_end - header_start 	                                ; Header length
    dd 0x100000000 - (0xe85250d6 + 0 + (header_end - header_start)) ; Checksum

;align 8
;framebuffer_tag_start:
;    dw 5                                              ; type
;    dw 1                                              ; flags
;    dd framebuffer_tag_end - framebuffer_tag_start    ; size
;    dd 800                                            ; width
;    dd 600                                            ; height
;    dd 32                                             ; depth
;framebuffer_tag_end:

align 8
    ; Required end tag:
    dw 0	; type
    dw 0	; flags
    dw 8	; size
header_end:

section .text
bits 32

_start:
    cli

    mov esp, stack_top

	push ebx
	push eax

    call main ; Jump main function

; hlt loop to avoid the Operating System immediately exiting
.hang: 
    hlt ; hlt instruction makes CPU wait for an interrupt before continuing
    jmp .hang

; Function is called from C with (uint32_t)&gdt_pointer as argument
gdt_flush:
	mov eax, [esp+4] ; Moves the GDT pointer from the stack to eax register
					; esp+4 is used to go past the return memory address 
					; pushed onto the stack and [] dereferences the 
					; value at that memory address
 
	lgdt [eax] ; Loads the GDT pointer into the CPU.

	mov ax, 0x10 ; Datasegment is at the GDT entry 0x10, and is loaded into the ax register
	mov ds, ax ; loads the address of the data segment GDT entry into all segment registers.
	mov es, ax
	mov fs, ax
	mov gs, ax
	mov ss, ax
	jmp 0x08:.flush ; Uses a jmp instruction to change the code segment to GDT entry 0x08,
					; .flush immediately returns.
.flush:
	ret


; idt_flush(uint32_t idt_ptr_address)
; Called from C with the address of our idt_ptr_t struct.
; The lidt instruction reads limit+base from that address and loads the IDT.
idt_flush:
    mov eax, [esp+4]   ; first argument: address of our idt_ptr_t
    lidt [eax]         ; load the IDT (same idea as lgdt for the GDT)
    ret

section .bss
stack_bottom:
    resb 4096 * 16
stack_top:
    \end{minted}
    \caption{Implementation of \texttt{multiboot2.asm}}
    \label{lst:multiboot2_asm}
\end{listing}


\begin{listing}[H]
    \begin{minted}[
        linenos,
        breaklines,
        frame=single,
        framesep=3mm,
        baselinestretch=1.1,
        numbersep=6pt,
        fontsize=\footnotesize,
        tabsize=4,
        bgcolor=white,
        style=friendly,
    ]{c}
#include "../include/gdt.h"
#include "../include/idt.h"
#include "../include/irq.h"
#include "../include/isr.h"
#include "../include/kernel/memory.h"
#include "../include/kernel/pit.h"
#include "../include/keyboard.h"
#include "../include/libc/stdbool.h"
#include "../include/libc/stddef.h"
#include "../include/libc/stdint.h"
#include "../include/libc/stdio.h"
#include "../include/matrix.h"
#include "../include/menu.h"
#include "../include/multiboot2.h"
#include "../include/piano.h"
#include "../include/program.h"
#include "lib/song/song.h"

extern uint32_t end; // The memory address of where the kernel ends

int kernel_main();

int main(uint32_t magic, struct multiboot_info *mb_info_addr) {
  // Requirements from assignment 2
  gdt_init();
  // Requirements from assignment 3
  idt_init();
  isr_init();
  irq_init();
  keyboard_init();
  // Requirements from assignment 4
  init_kernel_memory(&end);
  init_paging();
  print_memory_layout();
  init_pit();

  // Enable interrupts after we disabled them in _start in multiboot2.h
  asm volatile("sti");

  int reload = 250;
  int counter = reload;

  // Main kernel loop
  while (true) {
    int entry = kb_dequeue(&kb);
    if (entry != -1) {
      keyboard_handler(entry);
    }
    if (active_program == PROGRAM_PIANO) {
      // Uses a counter and reload value to make the matrix
      // update happen less frequently.
      if (--counter == 0) {
        counter = reload;
        matrix_rain_frame();
      }
    }
    sleep_interrupt(1);
  }

  return kernel_main();
}
    \end{minted}
    \caption{Implementation of \texttt{kernel.c}}
    \label{lst:kernel_c}
\end{listing}


\begin{listing}[H]
    \begin{minted}[
        linenos,
        breaklines,
        frame=single,
        framesep=3mm,
        baselinestretch=1.1,
        numbersep=6pt,
        fontsize=\footnotesize,
        tabsize=4,
        bgcolor=white,
        style=friendly,
    ]{c}
#include "../../include/libc/stdio.h"  //
#include "../../include/libc/stdarg.h" //
#include "../../include/libc/stddef.h"
#include "../../include/libc/stdint.h"

#define VGA_WIDTH 80       // Width of screen, number of columns
#define VGA_HEIGHT 25      // Height of screen, number of rows
#define VGA_MEMORY 0xB8000 // Memory location of VGA memory buffer

static uint16_t *terminal_buffer =
    (uint16_t *)VGA_MEMORY;  // Terminal buffer, each index represents one
                             // character and one colour.
static int terminal_row = 0; // used to index the terminal buffer
static int terminal_col = 0; // used to index the terminal buffer
static uint8_t terminal_color = 0x07; // light grey on black.

// terminal_scroll moves every line up by one and cleans the last row.
static void terminal_scroll(uint8_t terminal_color) {
  for (int row = 1; row < VGA_HEIGHT; row++) {
    for (int col = 0; col < VGA_WIDTH; col++) {
      terminal_buffer[(row - 1) * VGA_WIDTH + col] =
          terminal_buffer[row * VGA_WIDTH + col];
    }
  }
  for (int col = 0; col < VGA_WIDTH; col++) {
    terminal_buffer[(VGA_HEIGHT - 1) * VGA_WIDTH + col] =
        (terminal_color << 8) | ' ';
  }
  terminal_row = VGA_HEIGHT - 1;
}

// writes a single character to the VGA buffer.
// uses terminal_col and terminal_row for indexing and state.
void terminal_putchar(char c, uint8_t terminal_color) {
  if (c == '\\n') {
    terminal_col = 0;
    terminal_row++;
  } else {
    int index = terminal_row * VGA_WIDTH + terminal_col;
    terminal_buffer[index] = (terminal_color << 8) | c;
    terminal_col++;

    if (terminal_col >= VGA_WIDTH) {
      terminal_col = 0;
      terminal_row++;
    }
  }

  if (terminal_row >= VGA_HEIGHT) {
    terminal_scroll(terminal_color);
  }
}

// Uses terminal_putchar to write strings.
static void terminal_write(const char *str) {
  for (int i = 0; str[i] != '\\0'; i++) {
    terminal_putchar(str[i], terminal_color);
  }
}

// Helper function that formats integer values to text. and prints them to
// the screen
static void print_int(int value) {
  if (value < 0) {
    terminal_putchar('-', terminal_color);
    value = -value;
  }

  char buf[12];
  int i = 0;

  if (value == 0) {
    terminal_putchar('0', terminal_color);
    return;
  }

  while (value > 0) {
    // uses mod to get the last digit.
    // uses '0' to get starting index of numbers in ASCII which we use for our
    // lookup table
    buf[i++] = '0' + (value % 10);
    // divides by 10 so the number shift one place to the right.
    value /= 10;
  }

  while (i > 0) {
    terminal_putchar(buf[--i], terminal_color);
  }
}

// helper function for printing hex values
static void print_hex(uint32_t value) {
  char hex_chars[] = "0123456789ABCDEF";
  // always prints 0x
  terminal_write("0x");

  char buf[8];
  int i = 0;

  if (value == 0) {
    terminal_putchar('0', terminal_color);
    return;
  }

  while (value > 0) {
    // Gets the last 4 bits which is the same as one hex value.
    buf[i++] = hex_chars[value & 0xF];
    // shifts all bits down by four. discarding the four lower bits.
    // Preparing for next iteration of the while loop
    value >>= 4;
  }

  while (i > 0) {
    terminal_putchar(buf[--i], terminal_color);
  }
}

// Helper function for printf, so that we can have printf with and without
// default colour.
static int vprintf_color(const char *format, va_list args,
                         uint8_t terminal_color) {
  // loops through every character and selects the appropriate helper function
  // to use for printing.
  for (int i = 0; format[i] != '\\0'; i++) {
    if (format[i] == '%' && format[i + 1] != '\\0') {
      i++;
      switch (format[i]) {
      case 's':
        terminal_write(va_arg(args, const char *));
        break;
      case 'd':
        print_int(va_arg(args, int));
        break;
      case 'x':
        print_hex(va_arg(args, uint32_t));
        break;
      case 'c':
        terminal_putchar((char)va_arg(args, int), terminal_color);
        break;
      case '%':
        terminal_putchar('%', terminal_color);
        break;
      }
    } else {
      terminal_putchar(format[i], terminal_color);
    }
  }
}

// The main function this file exposes.
// Every other function is a helper function.
int printf(const char *format, ...) {
  va_list args;
  va_start(args, format);
  vprintf_color(format, args, terminal_color);
  va_end(args);
  return 0;
}

int printf_color(const char *format, uint8_t terminal_color, ...) {
  // Uses va_list, va_start, va_end to have a undefined amount of arguments
  // given to the function.
  va_list args;
  va_start(args, terminal_color);
  vprintf_color(format, args, terminal_color);
  va_end(args);
  return 0;
}

// set specific index to some specific character and colour
void terminal_set_char(int x, int y, char c, uint8_t terminal_color) {
  if (x < 0 || x >= VGA_WIDTH || y < 0 || y >= VGA_HEIGHT)
    return;
  terminal_buffer[y * VGA_WIDTH + x] = (terminal_color << 8) | c;
}

void clearTerminal() {
  for (int x = 0; x < VGA_WIDTH; x++) {
    for (int y = 0; y < VGA_HEIGHT; y++) {
      terminal_set_char(x, y, ' ', terminal_color);
    }
  }
  terminal_row = 0;
  terminal_col = 0;
}

void terminal_write_color(const char *str, uint8_t terminal_color) {
  for (int i = 0; str[i] != '\\0'; i++) {
    terminal_putchar(str[i], terminal_color);
  }
}
    \end{minted}
    \caption{Implementation of \texttt{stdio.c}}
    \label{lst:stdio_c}
\end{listing}


\subsection{Assignment 3}
\label{apx:assignment_3}


\subsubsection{IDT Implementation}
\label{apx:idt_implementation}


\begin{listing}[H]
    \begin{minted}[
        linenos,
        breaklines,
        frame=single,
        framesep=3mm,
        baselinestretch=1.1,
        numbersep=6pt,
        fontsize=\footnotesize,
        tabsize=4,
        bgcolor=white,
        style=friendly,
    ]{c}
#ifndef IDT_H
#define IDT_H

#include "libc/stdint.h"

// Struct defines an entry in the interrupt description table IDT
typedef struct {
  uint16_t offset_low; // Offset is the addres of the interrupt service routine
                       // relative to its segment.
  uint16_t selector;   // Selector is the GDT entry where the ISR resides.
                     // In our system the code segment lives at 0x08 in the GDT
  uint8_t zero;      // Reserved bit always zero
  uint8_t type_attr; // Gate type + DPL + present bit
                     // bit 7 present bit = 1, entry is valid
                     // bits 6-5 DPL = 00, highest privelege level
  // bit 4  = 0, unused value for gate typer other than Task
  // gate.
  // bits 3-0 gate type, = 1110, 32 bit interrupt type.
  uint16_t offset_high;
} __attribute__((packed)) idt_entry_t;

// The IDT pointer, this is loaded with the lidt assembly instruction.
typedef struct {
  uint16_t limit; // Size of the IDT in bytes, subtracted by 1
  uint32_t base;  // Memory address to the start of the IDT
} __attribute__((packed)) idt_ptr_t;

// See the idt.c file for more information.
void idt_set_entry(uint8_t index, uint32_t handler, uint16_t selector,
                   uint8_t type_attr);

// See the idt.c file for more information.
void idt_init(void);

// Defined in the multiboot2.asm assembly file. Loads the IDT pointer into the
// CPU registers.
extern void idt_flush(uint32_t idt_ptr_address);

#endif
    \end{minted}
    \caption{Implementation of \texttt{idt.h}}
    \label{lst:idt_h}
\end{listing}


\begin{listing}[H]
    \begin{minted}[
        linenos,
        breaklines,
        frame=single,
        framesep=3mm,
        baselinestretch=1.1,
        numbersep=6pt,
        fontsize=\footnotesize,
        tabsize=4,
        bgcolor=white,
        style=friendly,
    ]{c}
#include "../../include/idt.h"

static idt_entry_t idt[256];
static idt_ptr_t idt_pointer;

// Helper function to populate interrupt descriptor table, since the memory
// layout is awkward due to historical reasons.
void idt_set_entry(uint8_t index, uint32_t handler, uint16_t selector,
                   uint8_t type_attr) {
  idt[index].offset_low = handler & 0xFFFF;
  idt[index].selector = selector;
  idt[index].zero = 0;
  idt[index].type_attr = type_attr;
  idt[index].offset_high = (handler >> 16) & 0xFFFF;
}

extern int32_t dummy_isr;

// Initialises up the interrupt descriptor table and then loads it into the CPU
// registers.
void idt_init(void) {
  idt_pointer.limit = sizeof(idt) - 1;
  idt_pointer.base = (uint32_t)&idt;

  // Creates an IDT entry for each of the 256 possible entries.
  for (int i = 0; i < 256; i++) {
    // Initialises all IDT entries with dummy ISR that immediatly returns.
    // 0x08 selector which is our code segment.
    // 0x8E 32 bit interrupt gate.
    idt_set_entry(i, dummy_isr, 0x08, 0x8E);
  }

  idt_flush((uint32_t)&idt_pointer);
}
    \end{minted}
    \caption{Implementation of \texttt{idt.c}}
    \label{lst:idt_c}
\end{listing}


\subsubsection{ISR Implementation}
\label{apx:isr_implementation}


\begin{listing}[H]
    \begin{minted}[
        linenos,
        breaklines,
        frame=single,
        framesep=3mm,
        baselinestretch=1.1,
        numbersep=6pt,
        fontsize=\footnotesize,
        tabsize=4,
        bgcolor=white,
        style=friendly,
    ]{c}
#ifndef ISR_H
#define ISR_H

#include "libc/stdint.h"

// Snapshot of all registers at the moment an interrupt fires.
// The layout must match exactly what isr_common_stub pushes onto the stack.
typedef struct {
  uint32_t ds; // The currently used datasegment
  uint32_t edi, esi, ebp, esp, ebx, edx, ecx,
      eax; // the following variables are pushed onto the stack by the `pusha`
           // asm instruction
  uint32_t int_no, err_code; // these variabels are pushed before our ISR stub
                             // in the isr.asm file
  uint32_t eip, cs, eflags, useresp,
      ss; // CPU automatically pushes these variables upon interrupt
} registers_t;

// Inserts ISR entries into the IDT
void isr_init(void);

// Called by isr_common_stub. Currently only prints interrupt number.
void isr_handler(registers_t *regs);

// Makes entrypoints to various interrupt service routines defined in the
// isr.asm file available in C.
extern void isr0(void); // Division By Zero
extern void isr1(void); // Debug
extern void isr2(void); // Non-Maskable Interrupt

#endif
    \end{minted}
    \caption{Implementation of \texttt{isr.h}}
    \label{lst:isr_h}
\end{listing}


\begin{listing}[H]
    \begin{minted}[
        linenos,
        breaklines,
        frame=single,
        framesep=3mm,
        baselinestretch=1.1,
        numbersep=6pt,
        fontsize=\footnotesize,
        tabsize=4,
        bgcolor=white,
        style=friendly,
    ]{c}
#include "../../include/isr.h"
#include "../../include/idt.h"
#include "../../include/libc/stdint.h"
#include "../../include/libc/stdio.h"

// Names for the first 3 CPU exceptions we handle
static const char *exception_messages[] = {
    "Division By Zero",      // ISR 0
    "Debug",                 // ISR 1
    "Non-Maskable Interrupt" // ISR 2
};

// Insert the three ISR stubs into the IDT.
// 0x08     = code segment selector, entry 1 in GDT
// 0x8E     = 32-bit interrupt gate, kernel level privelege , present. See the
// idt_entry_t struct in idt.h for more info.
void isr_init(void) {
  idt_set_entry(0, (uint32_t)isr0, 0x08, 0x8E);
  idt_set_entry(1, (uint32_t)isr1, 0x08, 0x8E);
  idt_set_entry(2, (uint32_t)isr2, 0x08, 0x8E);
}

// Called from isr_common_stub with a pointer to the saved register state.
void isr_handler(registers_t *regs) {
  printf("ISR fired: interrupt %d: %s\\n", regs->int_no,
         exception_messages[regs->int_no]);
}
    \end{minted}
    \caption{Implementation of \texttt{isr.c}}
    \label{lst:isr_c}
\end{listing}


\begin{listing}[H]
    \begin{minted}[
        linenos,
        breaklines,
        frame=single,
        framesep=3mm,
        baselinestretch=1.1,
        numbersep=6pt,
        fontsize=\footnotesize,
        tabsize=4,
        bgcolor=white,
        style=friendly,
    ]{nasm}
extern isr_handler  ; the C function we call from the common stub

global dummy_isr
global isr0
global isr1
global isr2

; ISR that immediatly returns; used to initialise all IDT entries.
dummy_isr: 
	iretd
; ISR 0, Division By Zero, the CPU does not automatically push an error code.
isr0:
    cli
    push dword 0   ; dummy error code
    push dword 0   ; interrupt number
    jmp isr_common_stub

; ISR 1, Debug, the CPU does not automatically push an error code.
isr1:
    cli
    push dword 0   ; dummy error code
    push dword 1   ; interrupt number
    jmp isr_common_stub

; ISR 2, Non-Maskable Interrupt, Debug, the CPU does not automatically push an error code.
isr2:
    cli
    push dword 0   ; dummy error code
    push dword 2   ; interrupt number
    jmp isr_common_stub


; Stub used by alL ISR. Handles pushing CPU registers and datasegment in accordance with the struct registers_t defined in isr.h
; Then handles calling the isr_handler with registers_t argument, then restores the CPU state.
isr_common_stub:
    pusha               ; push eax, ecx, edx, ebx, esp, ebp, esi, edi

	; NOTE: currently our operating system only has a single datasegment, but to
	; allow for compatability with user space in the future we save the current datasegment, load the
	; kernel datasegment, and restore the datasegment after the ISR handle finishes.
    mov ax, ds          ; save the current data segment
    push eax            ; pushes the extended ax register to ensure 32 bit alignment

    mov ax, 0x10        ; load the kernel data segment into every segment register
    mov ds, ax
    mov es, ax
    mov fs, ax
    mov gs, ax

    push esp            ; Pushes the pointer to the top of the stack to the stack. which represents the register_t struct used by the isr_handler.
    call isr_handler    ; Call the C handler.
    add esp, 4          ; Moves the stack pointer forward with 4 bytes moving past the stack pointer we pushed to the stack.

    pop eax             ; Pops the data segment back off the stack. And restores it back into all segment registers
    mov ds, ax           
    mov es, ax
    mov fs, ax
    mov gs, ax

    popa                ; Restores the general purpose registers pushed onto the stack by `pusha`.
    add esp, 8          ; Move the stack pointer 8 bytes forward to discard the interrupt number and error code we pushed onto the stack. 
    sti
    iret                ; Return from interrupt, handles popping of the elements which the CPU automatically pushed to the stack upon interrupt.


    \end{minted}
    \caption{Implementation of \texttt{isr.asm}}
    \label{lst:isr_asm}
\end{listing}


\subsubsection{IRQ Implementation}
\label{apx:irq_implementation}


\begin{listing}[H]
    \begin{minted}[
        linenos,
        breaklines,
        frame=single,
        framesep=3mm,
        baselinestretch=1.1,
        numbersep=6pt,
        fontsize=\footnotesize,
        tabsize=4,
        bgcolor=white,
        style=friendly,
    ]{c}
#ifndef IRQ_H

// for in-depth explanations of the following variables look to
// https://wiki.osdev.org/8259_PIC
#define IRQ_H
#define PIC1 0x20            // address to the master PIC
#define PIC2 0xA0            // address to the slave PIC
#define PIC1_COMMAND PIC1    // address to the MASTER PIC command port
#define PIC1_DATA (PIC1 + 1) // address to the MASTER PIC data port
#define PIC2_COMMAND PIC2    // address to the SLAVE PIC command port
#define PIC2_DATA (PIC2 + 1) // address to the SLAVE PIC DATA port

#define ICW1_INIT 0x10 // Initialization sent to command port on startup
#define ICW1_ICW4                                                              \\
  0x01 // Used on command port indicates that ICW4 will be present

#define ICW4_8086 0x01 // 8086 CPU mode
#define CASCADE_IRQ 2

#include "isr.h" // reuse registers_t

// More info in the irq.c file
void irq_init(void);

// More info in the irq.c file
void irq_handler(registers_t *regs);

// More info in the irq.c file
void irq_register_handler(int irq, void (*handler)(registers_t *));

// Makes entrypoints to various interrupt request defined in the irq.asm file
// available in C.
extern void irq0(void);
extern void irq1(void);
extern void irq2(void);
extern void irq3(void);
extern void irq4(void);
extern void irq5(void);
extern void irq6(void);
extern void irq7(void);
extern void irq8(void);
extern void irq9(void);
extern void irq10(void);
extern void irq11(void);
extern void irq12(void);
extern void irq13(void);
extern void irq14(void);
extern void irq15(void);

#endif
    \end{minted}
    \caption{Implementation of \texttt{irq.h}}
    \label{lst:irq_h}
\end{listing}


\begin{listing}[H]
    \begin{minted}[
        linenos,
        breaklines,
        frame=single,
        framesep=3mm,
        baselinestretch=1.1,
        numbersep=6pt,
        fontsize=\footnotesize,
        tabsize=4,
        bgcolor=white,
        style=friendly,
    ]{c}
#include "../../include/irq.h"
#include "../../include/idt.h"
#include "../../include/io.h"

// PIC remapping
// The PIC has a has 16 irq lines spread over two chips, a master chip and a
// slave chip. Each of them have a port to send data to and commands to.
static void pic_remap(void) {
  outb(PIC1_COMMAND, ICW1_INIT | ICW1_ICW4); // ICW1: init master PIC
  outb(PIC2_COMMAND, ICW1_INIT | ICW1_ICW4); // ICW1: init slave PIC

  outb(PIC1_DATA,
       0x20); // ICW2: master offsets IRQ 0-7 by 32 bits to interrupts 32-39
  outb(PIC2_COMMAND,
       0x28); // ICW2: slave offsets  IRQ 8-15 by 40 bytes to interrupts 40-47

  outb(PIC1_DATA, 1 << CASCADE_IRQ); // ICW3: Bit mask tells master there is
                                     // exactly one slave on irq line 2
  outb(PIC2_DATA, CASCADE_IRQ); // ICW3: Tells slave which irq line it is on in
                                // relation to the master.

  outb(PIC1_DATA, ICW4_8086); // ICW4: 8086 mode on master
  outb(PIC2_DATA, ICW4_8086); // ICW4: 8086 mode on slave

  outb(PIC1_DATA, 0); // unmask all interrupts on master
  outb(PIC2_DATA, 0); // unmask all interrupts on slave
}

// Initialise interrupt requests.
void irq_init(void) {
  pic_remap();

  // Inserts the 16 stubs into IDT.
  // Code segment: 0x08
  // 0x8E Present,Kernel level privelege, interrupt gate.
  idt_set_entry(32, (uint32_t)irq0, 0x08, 0x8E);
  idt_set_entry(33, (uint32_t)irq1, 0x08, 0x8E);
  idt_set_entry(34, (uint32_t)irq2, 0x08, 0x8E);
  idt_set_entry(35, (uint32_t)irq3, 0x08, 0x8E);
  idt_set_entry(36, (uint32_t)irq4, 0x08, 0x8E);
  idt_set_entry(37, (uint32_t)irq5, 0x08, 0x8E);
  idt_set_entry(38, (uint32_t)irq6, 0x08, 0x8E);
  idt_set_entry(39, (uint32_t)irq7, 0x08, 0x8E);
  idt_set_entry(40, (uint32_t)irq8, 0x08, 0x8E);
  idt_set_entry(41, (uint32_t)irq9, 0x08, 0x8E);
  idt_set_entry(42, (uint32_t)irq10, 0x08, 0x8E);
  idt_set_entry(43, (uint32_t)irq11, 0x08, 0x8E);
  idt_set_entry(44, (uint32_t)irq12, 0x08, 0x8E);
  idt_set_entry(45, (uint32_t)irq13, 0x08, 0x8E);
  idt_set_entry(46, (uint32_t)irq14, 0x08, 0x8E);
  idt_set_entry(47, (uint32_t)irq15, 0x08, 0x8E);
}

// IRQ handler table
static void (*irq_handlers[16])(registers_t *) = {0};
// Function which adds IRQ handlers to the IRQ handler table
void irq_register_handler(int irq, void (*handler)(registers_t *)) {
  irq_handlers[irq] = handler;
}

// IRQ handler
void irq_handler(registers_t *regs) {
  int irq = regs->int_no - 32;

  if (irq_handlers[irq]) {
    irq_handlers[irq](regs);
  }
  // IRQ 0-7 are mapped to the master PIC. And IRQ 8-15 are mapped to the slave
  // PIC. Since we have shifted all interrupts up by 32 entries, that means that
  // the the slave PIC uses interrupt numbers starting from 40. Therefore we
  // must send EOI to both the slave and master PIC when the interrupt number is
  // larger or equal to 40.
  if (regs->int_no >= 40) {
    outb(PIC2_COMMAND, 0x20); // Send EOI to the slave PIC
  }
  outb(PIC1_COMMAND, 0x20); // Send EOI to the master PIC
}
    \end{minted}
    \caption{Implementation of \texttt{irq.c}}
    \label{lst:irq_c}
\end{listing}


\begin{listing}[H]
    \begin{minted}[
        linenos,
        breaklines,
        frame=single,
        framesep=3mm,
        baselinestretch=1.1,
        numbersep=6pt,
        fontsize=\footnotesize,
        tabsize=4,
        bgcolor=white,
        style=friendly,
    ]{nasm}
extern irq_handler ; C function that handles IRQ

global irq0 ; Defines entrypoints to the different IRQs
global irq1
global irq2
global irq3
global irq4
global irq5
global irq6
global irq7
global irq8
global irq9
global irq10
global irq11
global irq12
global irq13
global irq14
global irq15

; Stub used by alL IRQ. Handles pushing CPU registers and datasegment in accordance with the struct registers_t defined in isr.h
; Then handles calling the isr_handler with registers_t argument, then restores the CPU state.
irq0:  cli
    push dword 0
    push dword 32
    jmp irq_common_stub

irq1:  cli
    push dword 0
    push dword 33
    jmp irq_common_stub

irq2:  cli
    push dword 0
    push dword 34
    jmp irq_common_stub

irq3:  cli
    push dword 0
    push dword 35
    jmp irq_common_stub

irq4:  cli
    push dword 0
    push dword 36
    jmp irq_common_stub

irq5:  cli
    push dword 0
    push dword 37
    jmp irq_common_stub

irq6:  cli
    push dword 0
    push dword 38
    jmp irq_common_stub

irq7:  cli
    push dword 0
    push dword 39
    jmp irq_common_stub

irq8:  cli
    push dword 0
    push dword 40
    jmp irq_common_stub

irq9:  cli
    push dword 0
    push dword 41
    jmp irq_common_stub

irq10: cli
    push dword 0
    push dword 42
    jmp irq_common_stub

irq11: cli
    push dword 0
    push dword 43
    jmp irq_common_stub

irq12: cli
    push dword 0
    push dword 44
    jmp irq_common_stub

irq13: cli
    push dword 0
    push dword 45
    jmp irq_common_stub

irq14: cli
    push dword 0
    push dword 46
    jmp irq_common_stub

irq15: cli
    push dword 0
    push dword 47
    jmp irq_common_stub

; See isr_common_stub in isr.asm for comments as this stub follows the exact same principles the only difference being calling irq_handler instead of isr_handler
irq_common_stub:
    pusha

    mov ax, ds
    push eax

    mov ax, 0x10        ; kernel data segment
    mov ds, ax
    mov es, ax
    mov fs, ax
    mov gs, ax

    push esp            ; pointer to registers_t
    call irq_handler
    add esp, 4

    pop eax
    mov ds, ax
    mov es, ax
    mov fs, ax
    mov gs, ax

    popa
    add esp, 8          ; remove int_no and err_code
    sti
    iret
    \end{minted}
    \caption{Implementation of \texttt{irq.asm}}
    \label{lst:irq_asm}
\end{listing}


\subsubsection{Keyboard logger}
\label{apx:keyboard_logger}


\begin{listing}[H]
    \begin{minted}[
        linenos,
        breaklines,
        frame=single,
        framesep=3mm,
        baselinestretch=1.1,
        numbersep=6pt,
        fontsize=\footnotesize,
        tabsize=4,
        bgcolor=white,
        style=friendly,
    ]{c}
#ifndef KEYBOARD_H
#define KEYBOARD_H

#include "../include/libc/stdint.h"
#define KEYBOARD_DATA_PORT 0x60
#define KEYBOARD_BUFFER_SIZE 128

// Stores keyboard scancodes in circular queue to be processed by the main
// kernel loop.
typedef struct {
  int buffer[KEYBOARD_BUFFER_SIZE];
  int front;
  int back;
} KeyboardBuffer;

extern KeyboardBuffer kb;
extern const char scancode_ascii[128];

void kb_enqueue(KeyboardBuffer *kb, int entry);
int kb_dequeue(KeyboardBuffer *kb);
void keyboard_init(void);

int get_key(int scancode);
void keyboard_handler(int scancode);

#endif
    \end{minted}
    \caption{Implementation of \texttt{keyboard.h}}
    \label{lst:keyboard_h}
\end{listing}


\begin{listing}[H]
    \begin{minted}[
        linenos,
        breaklines,
        frame=single,
        framesep=3mm,
        baselinestretch=1.1,
        numbersep=6pt,
        fontsize=\footnotesize,
        tabsize=4,
        bgcolor=white,
        style=friendly,
    ]{c}
#include "../../include/keyboard.h"
#include "../../include/fedrelandet.h"
#include "../../include/io.h"
#include "../../include/irq.h"
#include "../../include/libc/stdint.h"
#include "../../include/libc/stdio.h"
#include "../../include/menu.h"
#include "../../include/piano.h"
#include "../../include/program.h"
#include "../../include/radio.h"
#include "../../include/shell.h"
#include "song/song.h"

// Maps scancode to character. keyboard Layout: US QWERTY
const char scancode_ascii[128] = {
    0,   0,    '1',  '2', '3',  '4', '5', '6', '7', '8', '9', '0', '-',
    '=', '\\b', '\\t', 'q', 'w',  'e', 'r', 't', 'y', 'u', 'i', 'o', 'p',
    '[', ']',  '\\n', 0,   'a',  's', 'd', 'f', 'g', 'h', 'j', 'k', 'l',
    ';', '\\'', '`',  0,   '\\\\', 'z', 'x', 'c', 'v', 'b', 'n', 'm', ',',
    '.', '/',  0,    '*', 0,    ' ', 0,   0,   0,   0,   0,   0,   0,
    0,   0,    0,    0,   0,    0,   '7', '8', '9', '-', '4', '5', '6',
    '+', '1',  '2',  '3', '0',  '.', 0,   0,   0,   0,   0,   0,   0,
    0,   0,    0,    0,   0,    0,   0,   0,   0,   0,   0,   0,   0,
    0,   0,    0,    0,   0,    0,   0,   0,   0,   0,   0,   0,   0,
    0,   0,    0,    0,   0,    0,   0,   0,   0,   0,   0,
};

// Function to add element to circular queue.
void kb_enqueue(KeyboardBuffer *kb, int entry) {
  kb->buffer[kb->front] = entry;
  kb->front = (kb->front + 1) % KEYBOARD_BUFFER_SIZE;
}
// Function to retrieve element to circular queue.
int kb_dequeue(KeyboardBuffer *kb) {
  if (kb->front == kb->back) {
    return -1;
  }
  int entry = kb->buffer[kb->back];
  kb->back = (kb->back + 1) % KEYBOARD_BUFFER_SIZE;
  return entry;
}

KeyboardBuffer kb = {0};

// Gets keyboard key, returning -1 on key release.
int get_key(int scancode) {
  if (scancode & 0x80)
    return -1;
  return scancode_ascii[scancode];
}

// Routes scancode to appropriate handler based on current active program.
void keyboard_handler(int scancode) {

  switch (active_program) {
  case PROGRAM_SHELL:
    shell_keyboard_handler(scancode);
    break;
  case PROGRAM_PIANO:
    piano_keyboard_handler(scancode);
    break;
  case PROGRAM_RADIO:
    radio_keyboard_handler(scancode);
    break;
  case PROGRAM_MENU:
    menu_keyboard_handler(scancode);
    break;
  case PROGRAM_FEDRELANDET:
    fedrelandet_keyboard_handler(scancode);
    break;
  default:
    printf("unknown program: %d\\n", active_program);
    break;
  }
}

// handler for IRQ1 (keyboard)
static void keyboard_irq_handler(registers_t *regs) {
  (void)regs;

  uint8_t scancode =
      inb(KEYBOARD_DATA_PORT); // Read scancode from keyboard data port

  // If 'q' is pressed set program to menu (quits current program)
  if (scancode == 0x10) {
    active_program = PROGRAM_MENU;
    disable_speaker();
    print_menu();
    return;
  }

  // queues keyboard scancode to ring buffer
  kb_enqueue(&kb, scancode);
}

// Maps IRQ 1 to the pointer of the keyboard handler function.
void keyboard_init(void) { irq_register_handler(1, keyboard_irq_handler); }
    \end{minted}
    \caption{Implementation of \texttt{keyboard.c}}
    \label{lst:keyboard_c}
\end{listing}


\subsection{Assignment 4}
\label{apx:assignment_4}


\subsubsection{Malloc}
\label{apx:malloc}


\begin{listing}[H]
    \begin{minted}[
        linenos,
        breaklines,
        frame=single,
        framesep=3mm,
        baselinestretch=1.1,
        numbersep=6pt,
        fontsize=\footnotesize,
        tabsize=4,
        bgcolor=white,
        style=friendly,
    ]{c}
#include "../../../include/kernel/memory.h"
#include "../../../include/libc/system.h"

#define MAX_PAGE_ALIGNED_ALLOCS 32

uint32_t last_alloc = 0;
uint32_t heap_end = 0;
uint32_t heap_begin = 0;
uint32_t pheap_begin = 0;
uint32_t pheap_end = 0;
uint8_t *pheap_desc = 0;
uint32_t memory_used = 0;

// Initialize the kernel memory manager
void init_kernel_memory(uint32_t* kernel_end)
{
    last_alloc = kernel_end + 0x1000;
    heap_begin = last_alloc;
    pheap_end = 0x400000;
    pheap_begin = pheap_end - (MAX_PAGE_ALIGNED_ALLOCS * 4096);
    heap_end = pheap_begin;
    memset((char *)heap_begin, 0, heap_end - heap_begin);
    pheap_desc = (uint8_t *)malloc(MAX_PAGE_ALIGNED_ALLOCS);
    printf("Kernel heap starts at 0x%x\\n", last_alloc);
}

// Print the current memory layout
void print_memory_layout()
{
    printf("Memory used: %d bytes\\n", memory_used);
    printf("Memory free: %d bytes\\n", heap_end - heap_begin - memory_used);
    printf("Heap size: %d bytes\\n", heap_end - heap_begin);
    printf("Heap start: 0x%x\\n", heap_begin);
    printf("Heap end: 0x%x\\n", heap_end);
    printf("PHeap start: 0x%x\\nPHeap end: 0x%x\\n", pheap_begin, pheap_end);
}

// Free a block of memory
void free(void *mem)
{
    alloc_t *alloc = (mem - sizeof(alloc_t));
    memory_used -= alloc->size + sizeof(alloc_t);
    alloc->status = 0;
}

// Free a block of page-aligned memory
void pfree(void *mem)
{
    if(mem < pheap_begin || mem > pheap_end) return;

    // Determine the page ID
    uint32_t ad = (uint32_t)mem;
    ad -= pheap_begin;
    ad /= 4096;

    // Set the page descriptor to free
    pheap_desc[ad] = 0;
}

// Allocate a block of page-aligned memory
char* pmalloc(size_t size)
{
    // Loop through the available list
    for(int i = 0; i < MAX_PAGE_ALIGNED_ALLOCS; i++)
    {
        if(pheap_desc[i]) continue;
        pheap_desc[i] = 1;
        printf("PAllocated from 0x%x to 0x%x\\n", pheap_begin + i*4096, pheap_begin + (i+1)*4096);
        return (char *)(pheap_begin + i*4096);
    }
    printf("pmalloc: FATAL: failure!\\n");
    return 0;
}


// Allocate a block of memory
void* malloc(size_t size)
{
    if(!size) return 0;

    // Loop through blocks to find an available block with enough size
    uint8_t *mem = (uint8_t *)heap_begin;
    while((uint32_t)mem < last_alloc)
    {
        alloc_t *a = (alloc_t *)mem;
        printf("mem=0x%x a={.status=%d, .size=%d}\\n", mem, a->status, a->size);

        if(!a->size)
            goto nalloc;
        if(a->status) {
            mem += a->size;
            mem += sizeof(alloc_t);
            mem += 4;
            continue;
        }
        // If the block is not allocated and its size is big enough,
        // adjust its size, set the status, and return the location.
        if(a->size >= size)
        {
            a->status = 1;
            printf("RE:Allocated %d bytes from 0x%x to 0x%x\\n", size, mem + sizeof(alloc_t), mem + sizeof(alloc_t) + size);
            memset(mem + sizeof(alloc_t), 0, size);
            memory_used += size + sizeof(alloc_t);
            return (char *)(mem + sizeof(alloc_t));
        }
        // If the block is not allocated and its size is not big enough,
        // add its size and the sizeof(alloc_t) to the pointer and continue.
        mem += a->size;
        mem += sizeof(alloc_t);
        mem += 4;
    }

    nalloc:;
    if(last_alloc + size + sizeof(alloc_t) >= heap_end)
    {
        panic("Cannot allocate bytes! Out of memory.\\n");
    }
    alloc_t *alloc = (alloc_t *)last_alloc;
    alloc->status = 1;
    alloc->size = size;

    last_alloc += size;
    last_alloc += sizeof(alloc_t);
    last_alloc += 4;
    printf("Allocated %d bytes from 0x%x to 0x%x\\n", size, (uint32_t)alloc + sizeof(alloc_t), last_alloc);
    memory_used += size + 4 + sizeof(alloc_t);
    memset((char *)((uint32_t)alloc + sizeof(alloc_t)), 0, size);
    return (char *)((uint32_t)alloc + sizeof(alloc_t));
}
    \end{minted}
    \caption{Implementation of \texttt{malloc.c}}
    \label{lst:malloc_c}
\end{listing}


\subsubsection{PIT Timer}
\label{apx:pit_timer}


\begin{listing}[H]
    \begin{minted}[
        linenos,
        breaklines,
        frame=single,
        framesep=3mm,
        baselinestretch=1.1,
        numbersep=6pt,
        fontsize=\footnotesize,
        tabsize=4,
        bgcolor=white,
        style=friendly,
    ]{c}
#include "../../include/kernel/pit.h"
#include "../../include/io.h"
#include "../../include/irq.h"
#include "../../include/libc/stdio.h"

// Counter for the amount of PIT ticks since init_pit() was called.
static volatile uint32_t tick = 0;

static uint32_t get_current_tick() { return tick; }

// Handler for IRQ0 called everytime the PIT triggers a hardware interrupt
static void pit_irq_handler(registers_t *regs) { tick++; }

// Initialise the PIT to fire at TARGET_FREQUENCY.
void init_pit() {
  irq_register_handler(0, pit_irq_handler);

  // Bit 7 - 6 = 00 = channel 0
  // Bit 5 - 4 = 11 = access mode: lobyte/hibyte
  // Bit 3 - 1 = 011 = square wave generator
  // bit 0 = 0 = 16-bit binary mode
  // 0011 0110 = 0x36
  uint16_t divisor = DIVIDER;
  outb(PIT_CMD_PORT, 0x36);
  outb(PIT_CHANNEL0_PORT, divisor & 0xFF);
  outb(PIT_CHANNEL0_PORT, (divisor >> 8) & 0xFF);
}

// Busy sleeping where the CPU runs useless while loop while waiting.
void sleep_busy(uint32_t milliseconds) {
  uint32_t start_tick = get_current_tick();
  uint32_t ticks_to_wait = milliseconds * TICKS_PER_MS;
  uint32_t elapsed_ticks = 0;

  while (elapsed_ticks < ticks_to_wait) {
    while (get_current_tick() == start_tick + elapsed_ticks) {
    }
    elapsed_ticks++;
  }
}

// Interrupt driven sleep where we allow interrupts with `sti` instruction,
// and put the CPU in a low resource usage mode waiting for hardware interrupt
// using the `hlt` instruction
void sleep_interrupt(uint32_t milliseconds) {
  uint32_t current_tick = get_current_tick();
  uint32_t ticks_to_wait = milliseconds * TICKS_PER_MS;
  uint32_t end_ticks = current_tick + ticks_to_wait;

  while (current_tick < end_ticks) {
    asm volatile("sti");
    asm volatile("hlt");
    current_tick = get_current_tick();
  }
}
    \end{minted}
    \caption{Implementation of \texttt{pit.c}}
    \label{lst:pit_c}
\end{listing}


\subsubsection{Song System Changes}
\label{apx:song_system_changes}


\begin{listing}[H]
    \begin{minted}[
        linenos,
        breaklines,
        frame=single,
        framesep=3mm,
        baselinestretch=1.1,
        numbersep=6pt,
        fontsize=\footnotesize,
        tabsize=4,
        bgcolor=white,
        style=friendly,
    ]{cpp}
#include "song/song.h"
#include "../../include/libc/stdint.h"
#include "../../include/program.h"

extern "C" {
#include "../../include/io.h"
#include "../../include/kernel/pit.h"
}

void enable_speaker() {
  // Read the current state of the PC speaker control register
  uint8_t speaker_state = inb(PC_SPEAKER_PORT);
  /*
  Bit 0: Speaker gate
          0: Speaker disabled
          1: Speaker enabled
  Bit 1: Speaker data
          0: Data is not passed to the speaker
          1: Data is passed to the speaker
  */
  // Check if bits 0 and 1 are not set (0 means that the speaker is disabled)
  if (speaker_state != (speaker_state | 3)) {
    // If bits 0 and 1 are not set, enable the speaker by setting bits 0 and 1
    // to 1
    outb(PC_SPEAKER_PORT, speaker_state | 3);
  }
}

extern "C" void disable_speaker() {
  // Turn off the PC speaker
  uint8_t speaker_state = inb(PC_SPEAKER_PORT);
  outb(PC_SPEAKER_PORT, speaker_state & 0xFC);
}

void play_sound(uint32_t frequency) {
  if (frequency == 0) {
    // If the note is no sound we disable the speaker
    disable_speaker();
    return;
  }

  auto divisor = (uint16_t)(PIT_BASE_FREQUENCY / frequency);

  // Set up the PIT
  outb(PIT_CMD_PORT, 0b10110110);
  outb(PIT_CHANNEL2_PORT, (uint8_t)(divisor & 0xFF));
  outb(PIT_CHANNEL2_PORT, (uint8_t)(divisor >> 8));
}

void play_song_impl(Song *song) {
  for (uint32_t i = 0; i < song->length; i++) {
    if (active_program != PROGRAM_RADIO &&
        active_program != PROGRAM_FEDRELANDET) {
      return;
    }
    // enable the speaker for each note in case previous note was no sound.
    enable_speaker();
    Note *note = &song->notes[i];
    // printf("Note: %d, Freq=%d, Sleep=%d\\n", i, note->frequency,
    // note->duration);
    play_sound(note->frequency);
    sleep_interrupt(note->duration);
  }
  disable_speaker();
}

void play_song(Song *song) { play_song_impl(song); }

extern "C" SongPlayer *create_song_player() {
  auto *player = new SongPlayer();
  player->play_song = play_song_impl;
  return player;
}

// enables speaker then plays sound. Used by piano.c
extern "C" void piano_play_sound(uint32_t freq) {
  enable_speaker();
  play_sound(freq);
}
    \end{minted}
    \caption{Implementation of \texttt{song.cpp}}
    \label{lst:song_cpp}
\end{listing}


\subsection{Assignment 5 }
\label{apx:assignment_5_}


\subsubsection{Menu System}
\label{apx:menu_system}


\begin{listing}[H]
    \begin{minted}[
        linenos,
        breaklines,
        frame=single,
        framesep=3mm,
        baselinestretch=1.1,
        numbersep=6pt,
        fontsize=\footnotesize,
        tabsize=4,
        bgcolor=white,
        style=friendly,
    ]{c}
#ifndef MENU_H
#define MENU_H

void print_menu();
void menu_navigate(int entry);
void menu_keyboard_handler(int scancode);
void print_radio_menu();

#endif
    \end{minted}
    \caption{Implementation of \texttt{menu.h}}
    \label{lst:menu_h}
\end{listing}


\begin{listing}[H]
    \begin{minted}[
        linenos,
        breaklines,
        frame=single,
        framesep=3mm,
        baselinestretch=1.1,
        numbersep=6pt,
        fontsize=\footnotesize,
        tabsize=4,
        bgcolor=white,
        style=friendly,
    ]{c}
#include "../../include/fedrelandet.h"
#include "../../include/keyboard.h"
#include "../../include/libc/stdio.h"
#include "../../include/program.h"

// Prints all available programs and their corresponding index to select them.
void print_menu() {
  printf("Select one of the following programs:\\n");
  for (int i = 0; i < PROGRAM_END; i++)
    printf("%d: %s\\n", i, program_names[i]);
}

void print_piano_menu() {
  printf("Welcome to the PIANO\\n");
  printf("To exit press 'q'\\n");
  printf("a = c\\n");
  printf("s = d\\n");
  printf("d = e\\n");
  printf("f = f\\n");
  printf("g = g\\n");
  printf("h = a\\n");
  printf("j = h\\n");
  printf("k = c\\n");
  printf("l = d\\n");
}

void print_radio_menu() {
  printf("Welcome to the RADIO\\n");
  printf("To exit press 'q'\\n");
  printf("press 0-5 to play a song :)\\n");
}

// switches active program based on key presses
void menu_navigate(int key) {
  // Translate ASCII numbers to actual numbers.
  int key_char = key - 48;
  switch (key_char) {
  case PROGRAM_SHELL:
    active_program = PROGRAM_SHELL;
    printf("ACTIVE PROGRAM: SHELL\\n");
    break;
  case PROGRAM_FEDRELANDET:
    active_program = PROGRAM_FEDRELANDET;
    printf("ACTIVE PROGRAM: FEDRELANDET\\n");
    hedre_fedrelandet();

    break;
  case PROGRAM_PIANO:
    active_program = PROGRAM_PIANO;
    clearTerminal();
    printf("ACTIVE PROGRAM: PIANO\\n");
    print_piano_menu();
    break;
  case PROGRAM_RADIO:
    active_program = PROGRAM_RADIO;
    print_radio_menu();
    printf("ACTIVE PROGRAM: RADIO\\n");
    break;
  default:
    active_program = PROGRAM_MENU;
    clearTerminal();
    printf("ACTIVE PROGRAM: MENU\\n");
    print_menu();
    break;
  }
}
void menu_keyboard_handler(int scancode) {
  if (scancode & 0x80)
    return;
  int key = scancode_ascii[scancode];
  menu_navigate(key);
}
    \end{minted}
    \caption{Implementation of \texttt{menu.c}}
    \label{lst:menu_c}
\end{listing}


\subsubsection{Program definitions}
\label{apx:program_definitions}


\begin{listing}[H]
    \begin{minted}[
        linenos,
        breaklines,
        frame=single,
        framesep=3mm,
        baselinestretch=1.1,
        numbersep=6pt,
        fontsize=\footnotesize,
        tabsize=4,
        bgcolor=white,
        style=friendly,
    ]{c}
#ifndef PROGRAM_H
#define PROGRAM_H

// Defines enums for all existing programs on KjebbeOS
typedef enum {
  PROGRAM_SHELL = 0,
  PROGRAM_PIANO,
  PROGRAM_RADIO,
  PROGRAM_MENU,
  PROGRAM_FEDRELANDET,
  // this must always be at the bottom
  PROGRAM_END,
} program_t;

static const char *program_names[] = {
    [PROGRAM_SHELL] = "shell",
    [PROGRAM_PIANO] = "piano",
    [PROGRAM_RADIO] = "radio",
    [PROGRAM_MENU] = "menu",
    [PROGRAM_FEDRELANDET] = "FEDRELANDET",
};

extern program_t active_program;

#endif
    \end{minted}
    \caption{Implementation of \texttt{program.h}}
    \label{lst:program_h}
\end{listing}


\begin{listing}[H]
    \begin{minted}[
        linenos,
        breaklines,
        frame=single,
        framesep=3mm,
        baselinestretch=1.1,
        numbersep=6pt,
        fontsize=\footnotesize,
        tabsize=4,
        bgcolor=white,
        style=friendly,
    ]{c}
#include "../../include/program.h"

// default program is menu
program_t active_program = PROGRAM_MENU;
    \end{minted}
    \caption{Implementation of \texttt{program.c}}
    \label{lst:program_c}
\end{listing}


\subsubsection{Keyboard Updates}
\label{apx:keyboard_updates}


\begin{listing}[H]
    \begin{minted}[
        linenos,
        breaklines,
        frame=single,
        framesep=3mm,
        baselinestretch=1.1,
        numbersep=6pt,
        fontsize=\footnotesize,
        tabsize=4,
        bgcolor=white,
        style=friendly,
    ]{c}
#ifndef KEYBOARD_H
#define KEYBOARD_H

#include "../include/libc/stdint.h"
#define KEYBOARD_DATA_PORT 0x60
#define KEYBOARD_BUFFER_SIZE 128

// Stores keyboard scancodes in circular queue to be processed by the main
// kernel loop.
typedef struct {
  int buffer[KEYBOARD_BUFFER_SIZE];
  int front;
  int back;
} KeyboardBuffer;

extern KeyboardBuffer kb;
extern const char scancode_ascii[128];

void kb_enqueue(KeyboardBuffer *kb, int entry);
int kb_dequeue(KeyboardBuffer *kb);
void keyboard_init(void);

int get_key(int scancode);
void keyboard_handler(int scancode);

#endif
    \end{minted}
    \caption{Implementation of \texttt{keyboard.h}}
    \label{lst:keyboard_h}
\end{listing}


\begin{listing}[H]
    \begin{minted}[
        linenos,
        breaklines,
        frame=single,
        framesep=3mm,
        baselinestretch=1.1,
        numbersep=6pt,
        fontsize=\footnotesize,
        tabsize=4,
        bgcolor=white,
        style=friendly,
    ]{c}
#include "../../include/keyboard.h"
#include "../../include/fedrelandet.h"
#include "../../include/io.h"
#include "../../include/irq.h"
#include "../../include/libc/stdint.h"
#include "../../include/libc/stdio.h"
#include "../../include/menu.h"
#include "../../include/piano.h"
#include "../../include/program.h"
#include "../../include/radio.h"
#include "../../include/shell.h"
#include "song/song.h"

// Maps scancode to character. keyboard Layout: US QWERTY
const char scancode_ascii[128] = {
    0,   0,    '1',  '2', '3',  '4', '5', '6', '7', '8', '9', '0', '-',
    '=', '\\b', '\\t', 'q', 'w',  'e', 'r', 't', 'y', 'u', 'i', 'o', 'p',
    '[', ']',  '\\n', 0,   'a',  's', 'd', 'f', 'g', 'h', 'j', 'k', 'l',
    ';', '\\'', '`',  0,   '\\\\', 'z', 'x', 'c', 'v', 'b', 'n', 'm', ',',
    '.', '/',  0,    '*', 0,    ' ', 0,   0,   0,   0,   0,   0,   0,
    0,   0,    0,    0,   0,    0,   '7', '8', '9', '-', '4', '5', '6',
    '+', '1',  '2',  '3', '0',  '.', 0,   0,   0,   0,   0,   0,   0,
    0,   0,    0,    0,   0,    0,   0,   0,   0,   0,   0,   0,   0,
    0,   0,    0,    0,   0,    0,   0,   0,   0,   0,   0,   0,   0,
    0,   0,    0,    0,   0,    0,   0,   0,   0,   0,   0,
};

// Function to add element to circular queue.
void kb_enqueue(KeyboardBuffer *kb, int entry) {
  kb->buffer[kb->front] = entry;
  kb->front = (kb->front + 1) % KEYBOARD_BUFFER_SIZE;
}
// Function to retrieve element to circular queue.
int kb_dequeue(KeyboardBuffer *kb) {
  if (kb->front == kb->back) {
    return -1;
  }
  int entry = kb->buffer[kb->back];
  kb->back = (kb->back + 1) % KEYBOARD_BUFFER_SIZE;
  return entry;
}

KeyboardBuffer kb = {0};

// Gets keyboard key, returning -1 on key release.
int get_key(int scancode) {
  if (scancode & 0x80)
    return -1;
  return scancode_ascii[scancode];
}

// Routes scancode to appropriate handler based on current active program.
void keyboard_handler(int scancode) {

  switch (active_program) {
  case PROGRAM_SHELL:
    shell_keyboard_handler(scancode);
    break;
  case PROGRAM_PIANO:
    piano_keyboard_handler(scancode);
    break;
  case PROGRAM_RADIO:
    radio_keyboard_handler(scancode);
    break;
  case PROGRAM_MENU:
    menu_keyboard_handler(scancode);
    break;
  case PROGRAM_FEDRELANDET:
    fedrelandet_keyboard_handler(scancode);
    break;
  default:
    printf("unknown program: %d\\n", active_program);
    break;
  }
}

// handler for IRQ1 (keyboard)
static void keyboard_irq_handler(registers_t *regs) {
  (void)regs;

  uint8_t scancode =
      inb(KEYBOARD_DATA_PORT); // Read scancode from keyboard data port

  // If 'q' is pressed set program to menu (quits current program)
  if (scancode == 0x10) {
    active_program = PROGRAM_MENU;
    disable_speaker();
    print_menu();
    return;
  }

  // queues keyboard scancode to ring buffer
  kb_enqueue(&kb, scancode);
}

// Maps IRQ 1 to the pointer of the keyboard handler function.
void keyboard_init(void) { irq_register_handler(1, keyboard_irq_handler); }
    \end{minted}
    \caption{Implementation of \texttt{keyboard.c}}
    \label{lst:keyboard_c}
\end{listing}


\begin{listing}[H]
    \begin{minted}[
        linenos,
        breaklines,
        frame=single,
        framesep=3mm,
        baselinestretch=1.1,
        numbersep=6pt,
        fontsize=\footnotesize,
        tabsize=4,
        bgcolor=white,
        style=friendly,
    ]{cpp}
#include "song/song.h"
#include "../../include/libc/stdint.h"
#include "../../include/program.h"

extern "C" {
#include "../../include/io.h"
#include "../../include/kernel/pit.h"
}

void enable_speaker() {
  // Read the current state of the PC speaker control register
  uint8_t speaker_state = inb(PC_SPEAKER_PORT);
  /*
  Bit 0: Speaker gate
          0: Speaker disabled
          1: Speaker enabled
  Bit 1: Speaker data
          0: Data is not passed to the speaker
          1: Data is passed to the speaker
  */
  // Check if bits 0 and 1 are not set (0 means that the speaker is disabled)
  if (speaker_state != (speaker_state | 3)) {
    // If bits 0 and 1 are not set, enable the speaker by setting bits 0 and 1
    // to 1
    outb(PC_SPEAKER_PORT, speaker_state | 3);
  }
}

extern "C" void disable_speaker() {
  // Turn off the PC speaker
  uint8_t speaker_state = inb(PC_SPEAKER_PORT);
  outb(PC_SPEAKER_PORT, speaker_state & 0xFC);
}

void play_sound(uint32_t frequency) {
  if (frequency == 0) {
    // If the note is no sound we disable the speaker
    disable_speaker();
    return;
  }

  auto divisor = (uint16_t)(PIT_BASE_FREQUENCY / frequency);

  // Set up the PIT
  outb(PIT_CMD_PORT, 0b10110110);
  outb(PIT_CHANNEL2_PORT, (uint8_t)(divisor & 0xFF));
  outb(PIT_CHANNEL2_PORT, (uint8_t)(divisor >> 8));
}

void play_song_impl(Song *song) {
  for (uint32_t i = 0; i < song->length; i++) {
    if (active_program != PROGRAM_RADIO &&
        active_program != PROGRAM_FEDRELANDET) {
      return;
    }
    // enable the speaker for each note in case previous note was no sound.
    enable_speaker();
    Note *note = &song->notes[i];
    // printf("Note: %d, Freq=%d, Sleep=%d\\n", i, note->frequency,
    // note->duration);
    play_sound(note->frequency);
    sleep_interrupt(note->duration);
  }
  disable_speaker();
}

void play_song(Song *song) { play_song_impl(song); }

extern "C" SongPlayer *create_song_player() {
  auto *player = new SongPlayer();
  player->play_song = play_song_impl;
  return player;
}

// enables speaker then plays sound. Used by piano.c
extern "C" void piano_play_sound(uint32_t freq) {
  enable_speaker();
  play_sound(freq);
}
    \end{minted}
    \caption{Implementation of \texttt{song.cpp}}
    \label{lst:song_cpp}
\end{listing}


\subsection{Assignment 6 }
\label{apx:assignment_6_}


\subsubsection{National Anthem }
\label{apx:national_anthem_}


\begin{listing}[H]
    \begin{minted}[
        linenos,
        breaklines,
        frame=single,
        framesep=3mm,
        baselinestretch=1.1,
        numbersep=6pt,
        fontsize=\footnotesize,
        tabsize=4,
        bgcolor=white,
        style=friendly,
    ]{c}
#ifndef FEDRELANDET_H
#define FEDRELANDET_H
void hedre_fedrelandet();
void fedrelandet_keyboard_handler(int scancode);

#endif
    \end{minted}
    \caption{Implementation of \texttt{fedrelandet.h}}
    \label{lst:fedrelandet_h}
\end{listing}


\begin{listing}[H]
    \begin{minted}[
        linenos,
        breaklines,
        frame=single,
        framesep=3mm,
        baselinestretch=1.1,
        numbersep=6pt,
        fontsize=\footnotesize,
        tabsize=4,
        bgcolor=white,
        style=friendly,
    ]{c}
#include "../../include/kernel/pit.h"
#include "../../include/keyboard.h"
#include "../../include/libc/stdio.h"
#include "../../include/matrix.h"
#include "./song/song.h"

Song nasjonal = {.notes = nasjonal_sangen,
                 .length =
                     sizeof(nasjonal_sangen) / sizeof(nasjonal_sangen[0])};

// Prints Norwegian flag and plays Norwegian national anthem
void hedre_fedrelandet() {
  SongPlayer *songPlayer = create_song_player();
  Song song_spaced = add_space_between_notes(&nasjonal, 50);
  clearTerminal();
  ja_vi_elsker();
  songPlayer->play_song(&song_spaced);
}

void fedrelandet_keyboard_handler(int scancode) {
  int key = get_key(scancode);
  if (key == 'j') {
    ja_vi_elsker();
  }
  if (key == 'h') {
    hedre_fedrelandet();
  }
}
    \end{minted}
    \caption{Implementation of \texttt{fedrelandet.c}}
    \label{lst:fedrelandet_c}
\end{listing}


\subsubsection{Radio Feature}
\label{apx:radio_feature}


\begin{listing}[H]
    \begin{minted}[
        linenos,
        breaklines,
        frame=single,
        framesep=3mm,
        baselinestretch=1.1,
        numbersep=6pt,
        fontsize=\footnotesize,
        tabsize=4,
        bgcolor=white,
        style=friendly,
    ]{c}
#ifndef RADIO_H
#define RADIO_H

void radio_keyboard_handler(int scancode);

#endif
    \end{minted}
    \caption{Implementation of \texttt{radio.h}}
    \label{lst:radio_h}
\end{listing}


\begin{listing}[H]
    \begin{minted}[
        linenos,
        breaklines,
        frame=single,
        framesep=3mm,
        baselinestretch=1.1,
        numbersep=6pt,
        fontsize=\footnotesize,
        tabsize=4,
        bgcolor=white,
        style=friendly,
    ]{c}
#include "../../include/radio.h"
#include "../../include/keyboard.h"
#include "../../include/libc/stdio.h"
#include "../../include/menu.h"
#include "../lib/song/song.h"

// Defines all songs available in song.h
Song song_1 = {.notes = music_1,
               .length = sizeof(music_1) / sizeof(music_1[0])};

Song song_2 = {.notes = music_2,
               .length = sizeof(music_2) / sizeof(music_2[0])};

Song song_3 = {.notes = music_3,
               .length = sizeof(music_3) / sizeof(music_3[0])};

Song song_4 = {.notes = music_4,
               .length = sizeof(music_4) / sizeof(music_4[0])};

Song song_5 = {.notes = music_5,
               .length = sizeof(music_5) / sizeof(music_5[0])};

Song song_6 = {.notes = music_6,
               .length = sizeof(music_6) / sizeof(music_6[0])};

Song song_7 = {.notes = nasjonal_sangen,
               .length = sizeof(nasjonal_sangen) / sizeof(nasjonal_sangen[0])};

Song *songs[] = {[0] = &song_1, [1] = &song_2, [2] = &song_3, [3] = &song_4,
                 [4] = &song_5, [5] = &song_6, [6] = &song_7};

void radio_keyboard_handler(int scancode) {
  int test_key = get_key(scancode);
  if (test_key == -1) {
    return;
  }
  int key = test_key - 48;

  SongPlayer *songPlayer = create_song_player();
  printf("Playing song number: %d", key);
  Song *song_to_play = songs[key];
  Song song_spaced = add_space_between_notes(song_to_play, 50);
  songPlayer->play_song(&song_spaced);
  printf("Finished playing song number: %d", key);
  print_radio_menu();
}
    \end{minted}
    \caption{Implementation of \texttt{radio.c}}
    \label{lst:radio_c}
\end{listing}


\subsubsection{Piano Application}
\label{apx:piano_application}


\begin{listing}[H]
    \begin{minted}[
        linenos,
        breaklines,
        frame=single,
        framesep=3mm,
        baselinestretch=1.1,
        numbersep=6pt,
        fontsize=\footnotesize,
        tabsize=4,
        bgcolor=white,
        style=friendly,
    ]{c}
#ifndef PIANO_H
#define PIANO_H

#include "libc/stdint.h"

void piano_keyboard_handler(uint8_t scancode);

#endif
    \end{minted}
    \caption{Implementation of \texttt{piano.h}}
    \label{lst:piano_h}
\end{listing}


\begin{listing}[H]
    \begin{minted}[
        linenos,
        breaklines,
        frame=single,
        framesep=3mm,
        baselinestretch=1.1,
        numbersep=6pt,
        fontsize=\footnotesize,
        tabsize=4,
        bgcolor=white,
        style=friendly,
    ]{c}
#include "../../include/piano.h"
#include "../../include/kernel/pit.h"
#include "../../include/keyboard.h"
#include "../../include/libc/stdio.h"
#include "../../include/matrix.h"
#include "song/frequencies.h"
#include "song/song.h"

// Defines piano keyboard layout and their offsets from C.
static const uint32_t scancode_note_offsets[128] = {
    [0x1E] = 1,  // a (c)
    [0x11] = 2,  // w (Db)
    [0x1F] = 3,  // s (d)
    [0x12] = 4,  // e (Eb)
    [0x20] = 5,  // d (e)
    [0x21] = 6,  // f (f)
    [0x14] = 7,  // t (f#)
    [0x22] = 8,  // g (g)
    [0x15] = 9,  // y (Ab)
    [0x23] = 10, // h (A)
    [0x16] = 11, // u (Bb)
    [0x24] = 12, // j (H)
    [0x25] = 13, // k (C)
    [0x18] = 14, // o (Db)
    [0x26] = 15, // l (d)
};

void piano_keyboard_handler(uint8_t scancode) {
  // if key release bit is set we disable speaker and return
  if (scancode & 0x80) {
    disable_speaker();
    return;
  }
  int init_offset = scancode_note_offsets[scancode];
  if (init_offset) {
    int note_offset = init_offset - 1;
    int octave = 5;
    int freq_to_play = (octave * 12) + note_offset;
    int freq = freqs[freq_to_play];
    // Changes the letter and number that is drawn on the matrix.
    update_matrix_current(note_offset);
    piano_play_sound(freq);
  }
}
    \end{minted}
    \caption{Implementation of \texttt{piano.c}}
    \label{lst:piano_c}
\end{listing}


\subsubsection{Matrix Application}
\label{apx:matrix_application}


\begin{listing}[H]
    \begin{minted}[
        linenos,
        breaklines,
        frame=single,
        framesep=3mm,
        baselinestretch=1.1,
        numbersep=6pt,
        fontsize=\footnotesize,
        tabsize=4,
        bgcolor=white,
        style=friendly,
    ]{c}
#ifndef MATRIX_H
#define MATRIX_H

void matrix_rain_frame();
void print_frame();
void colorshow();
void ja_vi_elsker();
void update_matrix_current();

extern int matrix_current_small;
extern int matrix_current_large;
extern int matrix_current_number;

#endif
    \end{minted}
    \caption{Implementation of \texttt{matrix.h}}
    \label{lst:matrix_h}
\end{listing}


\begin{listing}[H]
    \begin{minted}[
        linenos,
        breaklines,
        frame=single,
        framesep=3mm,
        baselinestretch=1.1,
        numbersep=6pt,
        fontsize=\footnotesize,
        tabsize=4,
        bgcolor=white,
        style=friendly,
    ]{c}
#include "../../include/kernel/pit.h"
#include "../../include/libc/stdio.h"

char *width_string = "0000000000000000000000000000000000000000000000000000000"
                     "0000000000000000000000000";

char *width1_string = "00000000000000000000";

char *width2_string = "00000";

char *width3_string = "0000000000";

char *width4_string = "00000";

char *width5_string = "0000000000000000000000000000000000000000";

// Defines top and bottom strips for the Norwegian flag
void print_strip() {
  printf_color(width1_string, VGA_RED);
  printf_color(width2_string, VGA_WHITE);
  printf_color(width3_string, VGA_BLUE);
  printf_color(width4_string, VGA_WHITE);
  printf_color(width5_string, VGA_RED);
}

// Prints top part of Norwegian flag
void print_top() {
  print_strip();
  print_strip();
  print_strip();
  print_strip();
  print_strip();
  print_strip();
  print_strip();
  print_strip();
}

// Prints bottom part of Norwegian flag
void print_bottom() {
  print_strip();
  print_strip();
  print_strip();
  print_strip();
  print_strip();
  print_strip();
  print_strip();
  print_strip();
}

void print_middle_strip() {
  printf_color(width1_string, VGA_WHITE);
  printf_color(width2_string, VGA_WHITE);
  printf_color(width3_string, VGA_BLUE);
  printf_color(width4_string, VGA_WHITE);
  printf_color(width5_string, VGA_WHITE);
}
// Prints middle part of Norwegian flag
void print_middle() {
  print_middle_strip();
  print_middle_strip();
  printf_color(width_string, VGA_BLUE);
  printf_color(width_string, VGA_BLUE);
  printf_color(width_string, VGA_BLUE);
  printf_color(width_string, VGA_BLUE);
  print_middle_strip();
  print_middle_strip();
}

// Main function to print Norwegian flag
void ja_vi_elsker() {
  clearTerminal();
  print_top();
  print_middle();
  print_bottom();
}

// pseudo random number generator
static uint32_t seed = 12345;
static uint32_t rng() {
  seed = seed * 1664525 + 1013904223;
  return seed;
}

int matrix_current_small = 'a';
int matrix_current_large = 'A';
int matrix_current_number = '0';

// Changes the characters to be printed based on piano offset
void update_matrix_current(int offset) {
  switch (offset % 12) {
  case (0):
    matrix_current_small = 'c';
    matrix_current_large = 'C';
    matrix_current_number = '0';
    break;
  case (1):
    matrix_current_small = 'c';
    matrix_current_large = 'C';
    matrix_current_number = '1';
    break;
  case (2):
    matrix_current_small = 'd';
    matrix_current_large = 'D';
    matrix_current_number = '2';
    break;
  case (3):
    matrix_current_small = 'd';
    matrix_current_large = 'D';
    matrix_current_number = '3';
    break;

  case (4):
    matrix_current_small = 'e';
    matrix_current_large = 'E';
    matrix_current_number = '4';
    break;

  case (5):
    matrix_current_small = 'f';
    matrix_current_large = 'F';
    matrix_current_number = '5';
    break;
  case (6):
    matrix_current_small = 'f';
    matrix_current_large = 'F';
    matrix_current_number = '6';
    break;
  case (7):
    matrix_current_small = 'g';
    matrix_current_large = 'G';
    matrix_current_number = '7';
    break;
  case (8):
    matrix_current_small = 'g';
    matrix_current_large = 'G';
    matrix_current_number = '8';
    break;
  case (9):
    matrix_current_small = 'a';
    matrix_current_large = 'A';
    matrix_current_number = '9';
    break;
  case (10):
    matrix_current_small = 'a';
    matrix_current_large = 'A';
    matrix_current_number = '0';
    break;
  case (11):
    matrix_current_small = 'b';
    matrix_current_large = 'B';
    matrix_current_number = '1';
    break;
  }
}

// Reload values for each character to be printed based on what the piano plays.
const int reload_small = 30;
int counter_small = reload_small;
const int reload_large = 25;
int counter_large = reload_large;
const int reload_number = 48;
int counter_number = reload_number;

// Renders a single frame of the matrix rain
void matrix_rain_frame() {
  static int drop[80] = {0};

  for (int x = 0; x < 80; x++) {
    int y = drop[x];

    // Piano coupled animation
    if (--counter_small == 0) {
      terminal_set_char(rng() % 80, rng() % 25, matrix_current_small,
                        VGA_RED_ON_BLACK);
      // Clear random index on screen.
      terminal_set_char(rng() % 80, rng() % 25, ' ', VGA_WHITE_ON_BLACK);
      counter_small = reload_small;
    }
    if (--counter_large == 0) {
      terminal_set_char(rng() % 80, rng() % 25, matrix_current_large,
                        VGA_BLUE_ON_BLACK);
      // Clear random index on screen.
      terminal_set_char(rng() % 80, rng() % 25, ' ', VGA_WHITE_ON_BLACK);
      counter_large = reload_large;
    }
    if (--counter_number == 0) {
      terminal_set_char((rng() * rng()) % 80, rng() % 25, matrix_current_number,
                        VGA_WHITE_ON_BLACK);
      // Clear random index on screen.
      terminal_set_char(rng() % 80, rng() % 25, ' ', VGA_WHITE_ON_BLACK);
      counter_number = reload_number;
    }

    // Matrix rainfall
    // Set lowest row to yellow colour
    terminal_set_char(x, y, '0' + (rng() % 10), VGA_YELLOW_ON_BLACK);
    // Set trail to green colour
    terminal_set_char(x, (y + 24) % 25, '0' + (rng() % 10),
                      VGA_LIGHT_GREEN_ON_BLACK);
    // Clear previous row
    terminal_set_char(x, (y + 23) % 25, ' ', VGA_WHITE_ON_BLACK);

    // 1/10 chance for raindrop to fall one step lower.
    if (rng() % 10 == 1) {
      drop[x] = (drop[x] + 1) % 25;
    }
  }
}
    \end{minted}
    \caption{Implementation of \texttt{matrix.c}}
    \label{lst:matrix_c}
\end{listing}


\subsubsection{Final kernel}
\label{apx:final_kernel}


\begin{listing}[H]
    \begin{minted}[
        linenos,
        breaklines,
        frame=single,
        framesep=3mm,
        baselinestretch=1.1,
        numbersep=6pt,
        fontsize=\footnotesize,
        tabsize=4,
        bgcolor=white,
        style=friendly,
    ]{c}
#include "../include/gdt.h"
#include "../include/idt.h"
#include "../include/irq.h"
#include "../include/isr.h"
#include "../include/kernel/memory.h"
#include "../include/kernel/pit.h"
#include "../include/keyboard.h"
#include "../include/libc/stdbool.h"
#include "../include/libc/stddef.h"
#include "../include/libc/stdint.h"
#include "../include/libc/stdio.h"
#include "../include/matrix.h"
#include "../include/menu.h"
#include "../include/multiboot2.h"
#include "../include/piano.h"
#include "../include/program.h"
#include "lib/song/song.h"

extern uint32_t end; // The memory address of where the kernel ends

int kernel_main();

int main(uint32_t magic, struct multiboot_info *mb_info_addr) {
  // Requirements from assignment 2
  gdt_init();
  // Requirements from assignment 3
  idt_init();
  isr_init();
  irq_init();
  keyboard_init();
  // Requirements from assignment 4
  init_kernel_memory(&end);
  init_paging();
  print_memory_layout();
  init_pit();

  // Enable interrupts after we disabled them in _start in multiboot2.h
  asm volatile("sti");

  int reload = 250;
  int counter = reload;

  // Main kernel loop
  while (true) {
    int entry = kb_dequeue(&kb);
    if (entry != -1) {
      keyboard_handler(entry);
    }
    if (active_program == PROGRAM_PIANO) {
      // Uses a counter and reload value to make the matrix
      // update happen less frequently.
      if (--counter == 0) {
        counter = reload;
        matrix_rain_frame();
      }
    }
    sleep_interrupt(1);
  }

  return kernel_main();
}
    \end{minted}
    \caption{Implementation of \texttt{kernel.c}}
    \label{lst:kernel_c}
\end{listing}

\end{document}
